\chapter*{Lời nói đầu}

\thispagestyle{empty}

\markboth{{\it Lời nói đầu}}{{\it Lời nói đầu}}
\addcontentsline{toc}{section}{{\bf  Lời nói đầu}\rm}

\vspace*{-0.2cm}

Sau khi một loạt cuốn sách về phương pháp giải toán được bạn đọc đón nhận []-[], những cuốn sách này liên tục được tái bản và nhiều bạn đọc khen hơn là chê. Điều đó động viên tôi thực hiện biên tập cuốn sách này, đúng như tên của cuốn sách là tuyển tập các phương pháp và các chuyên đề giải toán chứ không phải tuyển tập các bài toán hay. Ta đã biết rất nhiều phương pháp hay đã được tôi biên tập trong các cuốn []-[], sau một thời gian tìm hiểu kĩ hơn nữa thì tôi thấy các phương pháp này giải được rất nhiều dạng bài toán khác nhau, trong tay tôi đã có rất nhiều tài liệu mà những cuốn sách trước không có được. Tôi biên tập cuốn sách này để củng cố các phương pháp giải toán mà các cuốn sách trước đã thể hiện và đưa thêm một số phương pháp khác, cách nhìn khác về việc giải toán. Đọc tài liệu này các bạn sẽ thấy tuy là phương pháp giải toán khác nhau nhưng nó có một tư tưởng thống nhất là {\it suy luận có lí}. Số bài tập hay dùng các phương pháp giải khác nhau là vô cùng nhiều, nên tất cả những bài toán trong các cuốn trước đây tôi không đưa vào đây. Tôi cố gắng chọn những bài toán hay, mới vào tuyển tập này. Nếu có những bài toán trùng với các tập sách trước thì sẽ có một cách giải hoàn toàn mới, bạn đọc có thể so sánh với những cách giải cũ. Cuốn sách được chia làm hai phần lớn:

Phần I. Các phương pháp giải toán.

1. Phương pháp chứng minh bằng phản chứng.

2. Phương pháp dùng ví dụ, phản ví dụ và xây dựng lời giải.

3. Phương pháp nguyên lí Đirichle

4. Phương pháp quy nạp toán học 

5. Phương pháp dùng đại lượng bất biến

6. Phương pháp dùng đại lượng cực biên

7. Phương pháp tô màu

8. Các phương pháp khác

Phần II. Những chuyên đề cơ bản

1. Tổ hợp rời rạc

2. Lí thuyết số

3. Bất đẳng thức

4. Dãy số

5. Đa thức

6. Phương trình hàm

7. Hình học

8. Thuật toán và trò chơi

Mỗi phần trên đều được triển khai từ dễ đến khó và một lôgic có lí cao. Các bài tập và ví dụ được giải cẩn thận và dễ hiểu nhất. Bạn đọc có thể tìm thấy những lời giải khác hay hơn, ngắn hơn nhưng nhằm mục đích mô tả phương pháp giải toán nên ở đây có thể dài hơn. Phần cuối của mỗi chương là lời giải ngay các bài tập trong chương đó, đánh số các ví dụ, bài tập là lần lượt cùng nhau cho đến hết chương.
 
Cuốn sách dành cho học sinh phổ thông yêu toán, học sinh khá giỏi môn toán, 
các thầy cô giáo, sinh viên đại học ngành toán, ngành tin học 
và những người yêu thích toán học phổ thông. 
Trong biên soạn không thể tránh khỏi sai sót và nhầm lẫn mong bạn đọc cho ý kiến.
 Mọi góp ý gửi về địa chỉ: Ban biên tập sách Toán, Nhà xuất bản Giáo dục, $187^{\mbox{b}}$ Giảng Võ, Hà Nội.
 
Tác giả cảm ơn ban biên tập Toán - Nhà xuất bản Giáo dục Hà Nội đã hết sức giúp đỡ để cuốn sách được in ra. 
 
\hfill Hà Nội, ngày 2 tháng 11 năm  2006

\hfill {\bf Nguyễn Hữu Điển}\hspace{0.5cm}
